\documentclass[10pt,a4paper]{amsart}
\usepackage[margin=19mm]{geometry}
\usepackage{amsmath,amssymb,mathtools}
\usepackage[scr=rsfs,cal=euler]{mathalfa}
\usepackage{xfrac}
\usepackage[unicode,hidelinks,bookmarks=false]{hyperref}
\newcommand{\ie}{i.e.\ }
\newcommand{\cf}{cf.\ }
\newcommand{\resp}{respectively\ }
\newcommand{\Subsection}{\subsection}
\providecommand{\nobreakdash}{\nobreak\hbox{-}}
\setlength{\emergencystretch}{2em}
\multlinegap0pt
\usepackage{amssymb}
\usepackage{enumerate}
\usepackage[shortlabels]{enumitem}
\usepackage{tikz}
\newtheorem{theorem}{Theorem}
\title{Matchings avoiding ordered patterns}
\author{J\'anos Bar\'at, Andrea Freschi, G\'eza T\'oth}
\date{Source: arXiv:2512.15461v2 (CC BY 4.0)}
\begin{document}
\maketitle

\noindent\emph{Source and licence.} Matchings avoiding ordered patterns. Version arXiv:2512.15461v2 (CC BY 4.0), distributed by arXiv under the Creative Commons Attribution 4.0 International licence, http://creativecommons.org/licenses/by/4.0/, as recorded in the arXiv licence metadata for this exact version and shown on the abstract page https://arxiv.org/abs/2512.15461v2. Full licence notice: \texttt{licenses/arxiv-2512.15461-CC-BY-4.0.txt}.\par\smallskip

\noindent\textit{Editorial scope.} Counted unit: the article's theorem thm:RamseyAltPaths (source0.tex lines 1047--1050). The article's complete original proof of it is delivered with it (source0.tex lines 1052--1097, one \texttt{proof} environment of 46 lines). No mathematics is written, repaired or re-derived: the statement and the proof are the article's own lines, in the article's own notation, and no retained line is substituted or re-flowed.  No further source-local context is needed: every cross-reference of every retained line resolves inside the two delivered spans.  Nothing was omitted from any retained span, so the package carries no omission record.  No retained line contains a citation, so the wrapper carries no bibliography and no bibliography entry.  No retained line contains a figure environment, an image-inclusion command or drawing code, so no image is delivered and the package declares no asset dependency.  Editorial formatting only, in addition: an AMS \texttt{amsart} preamble that restates the article's own declarations and macros used by the retained lines, namely the environments \texttt{theorem} on separate counters, together with the standard packages those lines need.  The retained environments print in this document as Theorem 1 in order of appearance, whereas the article prints the same items with the numbering of its own full document.  The complete native source of the article as distributed in the arXiv e-print for this exact version is bundled unchanged as \texttt{sources/191176/source0.tex}.
\begin{theorem}\label{thm:RamseyAltPaths}
For every $t\in\mathbb N$, we have
$$R_<(\text{Alt-}P_{t})\le3t+3.$$    
\end{theorem}

\begin{proof}
    Note that it suffices to show that $R_<(\text{Alt-}P_{t})\le3t$ if $t$ is even.
    Indeed, if $t$ is odd, it follows that $R_<(\text{Alt-}P_{t})\le R_<(\text{Alt-}P_{t+1})\le3(t+1)=3t+3$.

    Let $t=2k$ for some $k\in\mathbb N$ and let $m:=R_<(\text{Alt-}P_{2k})-1$.
    By definition of $m$, there exists a red/blue edge-coloured complete graph $K_m$ with vertex set $[m]$ which does not contain a monochromatic copy of $\text{Alt-}P_{2k}$.
    Note that $R_<(\text{Alt-}P_{2k})\ge|V(\text{Alt-}P_{2k})|=2k$ and so $m\ge2k-1$.

    Let $A:=\{xy\in E(K_m):x+y\le 2k-2\}$ and $B:=\{xy\in E(K_m):x+y\ge 2m-2k+4\}$.
    Observe that 
    \begin{align*}
    |A|=|\{xy\in E(K_m):x+y\le 2k-2\}| &=(2k-4)+(2k-6)+\dots+2\\
    & =2[(k-2)+(k-3)+\dots+1]=2\binom{k-1}{2}
    \end{align*}
    and similarly
    $$|B|=|\{xy\in E(K_m):x+y\ge 2m-2k+4\}|=2\binom{k-1}{2}.$$
    Furthermore, $A\cap B=\emptyset$.
    Indeed, we cannot have $2m-2k+4\le x+y\le 2k-2$ since $m\ge2k-1$.

    It follows that
    $$|E(K_m)\setminus(A\cup B)|=\binom{m}{2}-|A|-|B|=\binom{m}{2}-4\binom{k-1}{2}.$$
    By the pigeonhole principle, at least half of the edges in $E(K_m)\setminus(A\cup B)$ have the same color, say red.

    The key observation is that edges in $A\cup B$ can be arbitrarily recolored without creating a monochromatic copy of $\text{Alt-}P_{2k}$.
    Indeed, suppose there exists a copy of $\text{Alt-}P_{2k}$ in the complete uncolored graph with vertex set $[m]$, which contains some edge $xy\in A$ with $x<y$.
    Since a copy of $\text{Alt-}P_{2k}$ can be decomposed into two edge-disjoint nested matchings of $k$ and $k-1$ edges respectively, we conclude that there exists a nested $(k-1)$-matching $M$ using the edge $xy$.
    Now, every edge of $M$, which lies above $xy$ must intersect $[x-1]$, while $xy$ and every edge of $M$ below $xy$ must lie in $[x,y]$.
    Therefore, we have
    $$e(M)\le|[x-1]|+\frac{|[x,y]|}{2}\le(x-1)+\frac{y-x+1}{2}=\frac{y+x-1}{2}<k-1$$
    where the last inequality follows from the assumption that $xy\in A$ and so $x+y\le 2k-2$.
    Hence $M$ has strictly less than $k-1$ edges, a contradiction.
    By symmetry, an analogous argument works for $xy\in B$.

    We conclude that we can recolor the edges in $A\cup B$ with red without creating a monochromatic alternating $2k$-path.
    Let $G$ be the graph with vertex set $[m]$ whose edges are the red edges in $K_m$.
    Hence, we have
    $$e(G)=|A|+|B|+\frac{1}{2}\left(\binom{m}{2}-|A|-|B|\right)=\frac{1}{2}\left(\binom{m}{2}+4\binom{k-1}{2}\right).$$
    Note that $G$ does not contain a copy of an alternating $2k$-path, as that would be a red monochromatic copy in $K_m$.
    Therefore, we have
    $$e(G)\le ex_<(n,\text{Alt-}P_{2k})=2(k-1)m-(k-1)(2k-1).$$

    Hence 
    $$\frac{1}{2}\left(\binom{m}{2}+4\binom{k-1}{2}\right)\le 2(k-1)m-(k-1)(2k-1)$$
    which implies $m\le6k=3t$.
    That is, $R_<(\text{Alt-}P_{t})\le3t$ for $t$ even, as required.
\end{proof}

\end{document}
